\section{Related Work}

\paragraph{Neural and LLM-guided theorem proving.}
Early generative provers demonstrated that language models can synthesize
machine-checkable proofs \cite{polu2020gptf}. In Lean, LeanDojo and ReProver
made retrieval-augmented proof search reproducible at scale
\cite{yang2023leandojo}, while Lean Copilot integrated learned suggestions,
goal completion, and premise selection into Lean workflows
\cite{song2024leancopilot}. Large-scale systems such as DeepSeek-Prover and its
successors improve proving through synthetic data, reinforcement learning,
subgoal decomposition, and high test-time sampling
\cite{xin2024deepseekprover,ren2025deepseekprov_v2}. Goedel-Prover and
Goedel-Prover-V2 similarly combine data synthesis and verifier-guided
self-correction \cite{lin2025goedel,lin2025goedelv2}. These works establish the
capability of the underlying provers. Our focus is complementary: how to
allocate a fixed inference budget among heterogeneous proof actions and
branches.

\paragraph{Verifier-guided agents and decomposition.}
Compiler feedback can make iterative proof generation substantially more
sample-efficient than independent generation. APOLLO combines Lean feedback,
proof repair, subgoal isolation, and automated solvers under a limited sampling
budget \cite{ospanov2025apollo}. A Minimal Agent shows that iterative
refinement, compact context management, and library search form a strong
baseline even without an elaborate agent hierarchy
\cite{requena2026minimal}. Compile to Compress uses structured compiler failure
modes for local refinement and value-guided search
\cite{li2026compile}. Goedel-Architect constructs and refines dependency
blueprints of definitions and lemmas \cite{chung2026architect}. These systems
motivate our use of verifier feedback and explicit dependencies. \system
differs in making the resource-allocation policy itself the object of study and
in keeping decomposition, repair, capability tests, representation changes,
and model escalation inside one typed action space.

\paragraph{Proof progress and cost-aware routing.}
HyperTree Proof Search uses learned search to improve neural theorem proving
\cite{lample2022hypertree}, and LeanProgress studies proof-progress prediction
as a search signal \cite{george2026leanprogress}. The closest concurrent work
routes an agent between continuing a target and restarting from a new
decomposition using predicted success and cost
\cite{rognvaldsson2026costquality}. Our controller generalizes this control
problem along three axes: a multi-resource rather than single compute budget, a
larger set of local and structural actions, and an explicit AND/OR workspace
from which downstream unlock value and partial completion can be measured.

\paragraph{Budget-aware test-time search.}
Outside theorem proving, Budget-Aware Value Tree Search conditions node
selection on remaining resources and uses residual rather than absolute value
estimation \cite{li2026bavt}. Budget-Guided MCTS changes its search behavior as
a fixed token budget is depleted \cite{miyamoto2026bgmcts}. We adapt these
ideas to a setting with a symbolic verifier. Lean feedback supplies typed
failure evidence, stable goal fingerprints, accepted intermediate facts, and
exact terminal correctness; the controller can therefore ground value in
state transitions that are unavailable in open-ended reasoning tasks.

\paragraph{Benchmarks.}
miniF2F provides cross-system Olympiad-style formal statements
\cite{zheng2021minif2f}; ProofNet pairs natural-language mathematics with Lean
statements and proofs \cite{azerbayev2023proofnet}; and PutnamBench targets
substantially harder undergraduate mathematics
\cite{tsoukalas2024putnambench}. We use fixed formal statements for the primary
controller comparison so that autoformalization errors do not confound search
allocation, then report a separate end-to-end setting with natural-language
inputs.
